\section{Related Work}
\label{sec:related}

This section situates the LDMs within the broader literature on generative models, statistical optimisation, and their intersection for scientific discovery. The aim is not to provide an exhaustive survey, but a focused positioning of LDM within the key strands of related work that directly motivate its design and delineate its contributions. We structure the discussion in three parts. We first review the remarkable progress in LLMs, including their scaling laws, emergent reasoning capabilities, and post-training methods; we then highlight the limitations of pure generative models for discovery, particularly the lack of calibrated uncertainty and expensive, sparse reward signals. Next, we survey the statistical learning paradigm for scientific discovery, from multi-armed bandits to Bayesian optimisation (BO), and discuss its strengths and challenges in high-dimensional or open-ended, unrepresented scientific spaces. Finally, we examine the growing body of work on hybrid approaches that combine LLMs with BO, clarify the differences between LDM and the most closely related methods, and articulate the open problems our framework addresses.

\subsection{Large Language Models and Their Limitations in Scientific Discovery}

The early development of language models was marked by several landmark architectures, including the Transformer \citep{vaswaniAttentionAllYou2017}, GPT-1 \citep{radfordImprovingLanguageUnderstanding2018}, and BERT \citep{devlinBERTPretrainingDeep2019}. The past few years have witnessed transformative progress in large language models. Scaling laws \citep{kaplanScalingLawsNeural2020} have established a predictable correlation in LLMs' performance with the scale of model weights, data, and test-time compute. The remarkable success of GPT-3 \citep{brownLanguageModelsAre2020} further propelled research along the scaling direction. As models scale, they exhibit emergent abilities \citep{weiEmergentAbilitiesLarge2022} not present in smaller counterparts, including complex reasoning via chain-of-thought prompting \citep{weiChainofThoughtPromptingElicits2023,kojimaLargeLanguageModels2023}, in-context learning \citep{hanUnderstandingEmergentInContext2025}, and instruction following \citep{ouyangTrainingLanguageModels2022}. In addition, post-training methods, such as supervised fine-tuning (SFT) and reinforcement learning from human feedback (RLHF) \citep{ouyangTrainingLanguageModels2022}, can better adapt LLMs to downstream tasks and the users' preferences.

Since 2023, \emph{Test-time compute} has been a crucial development that pushes LLMs from language understanding to reasoners for complex problems, where the key insight is that the use of additional computation at inference time greatly improves output quality. Early research utilises verifiers \citep{cobbeTrainingVerifiersSolve2021, lightmanLetsVerifyStep2023} to guide reasoning to solve complex mathematical problems. Further, \emph{tree of thoughts} \citep{yao2023tree} performs deliberate search over intermediate reasoning states with self-evaluation. Self-Refine \citep{madaan2023selfrefine} iterates generate--critique--revise cycles. More generally, \citet{snell2025scaling} show that optimal allocation of test-time compute allows small models to outperform much larger models. The underlying search machinery, including UCT \citep{kocsis2006bandit} and the PUCT rule popularised by AlphaZero \citep{silver2017mastering}, has been applied to LLM decoding and training \citep{fengAlphazerolikeTreeSearchCan2024}. The combination of inference-time search with verifier-guided refinement underpins renowned systems such as OpenAI's o1 \citep{openaiOpenAIO1System2024} and DeepSeek-R1 \citep{deepseek-aiDeepSeekR1IncentivizingReasoning2025}. When cheap supervision is repeatedly available, \emph{reinforcement learning with verifiable rewards} (RLVR) \citep{deepseek-aiDeepSeekR1IncentivizingReasoning2025, shaoDeepSeekMathPushingLimits2024} has led to impressive results in mathematics, coding, and reasoning tasks. A unified tutorial perspective on these methods is provided by \citet{wangTutorialLLMReasoning2025}, and OpenR \citep{wang2024openr} further establishes a comprehensive technical framework.

Despite the significant progress of LLM reasoning, how LLMs engage in open-ended scientific discovery remains questionable. Studies show that LLMs are also poorly calibrated for scientific decisions---their internal confidence does not reliably reflect uncertainty about an external property \citep{guptaLLMsBayesianOptimization2025,kristiadiSoberLookLLMs2024}. Moreover, end-to-end execution benchmarks consistently show that frontier LLM agents cannot yet reliably discover novel, high-quality solutions: on ResearchClawBench the strongest autonomous agent scores only 21.5/50 across 40 real-paper-derived tasks \citep{xu2026researchclawbench}, and on PaperBench agents replicate ICML papers at 21.0\% against 41.4\% by PhD researchers \citep{starace2025paperbench}. Even at the ideation stage, novelty rankings that initially favour LLMs \citep{si2024novelideas} reverse once ideas are executed \citep{si2025executiongap}. LLM-generated ideas also exhibit narrow exploration, converging on the centroid of existing literature \citep{tang2026narrowexploration}. 

Therefore, scientific discovery essentially requires a novel paradigm of foundation models.  Crucially, the key challenge is not intrinsic LLM failures: when proposals are paired with automated evaluators and search, genuinely novel results are recovered \citep{novikov2025alphaevolve}. The bottleneck is the absence of a calibrated external value signal.  Test-time compute and reinforcement learning typically rely on cheap verifiers or reward models, such as ground-truth answers or trained networks from large data, to guide the search of reasoning steps. This is built upon access to rich feedback data or ground-truth answers, which, however, is usually scarce in scientific domains.

\subsection{Statistical Learning for Scientific Discovery}

A mature body of statistical learning approaches addresses the problem of optimising a black-box reward function with noisy evaluations. \emph{Multi-arm bandit} constitutes a representative framework conceptualising the exploration--exploitation trade-off \citep{lattimoreBanditAlgorithms2020}, where the upper confidence bound (UCB) \citep{auerUsingConfidenceBounds2002} provides a simple way to balance these competing objectives with sublinear regret guarantees. However, classical multi-arm bandit settings assume a finite set of arms, yet scientific design spaces are often effectively unbounded.

This challenge of unbounded space is then captured by the \emph{infinitely-armed bandit} \citep{berry1997bandit, wang2008infinitely} frameworks, where new arms must be discovered from a \emph{reservoir}. The performance in such settings is governed by the reservoir's tail of near-optimal arms, a quantity known as the \emph{discovery gap}. The LDM framework explicitly adopts this view, treating the LLM as an adaptive reservoir and analysing the discovery gap through a missing-mass argument \citep{mcallester2000convergence, berend2013concentration} in \S~\ref{sec:theory-main}. In comparison, LDM utilises an LLM as its reservoir, endowing the search process with broad scientific priors and rich semantic knowledge to inform candidate generation. In addition, the LLM’s context $\mathcal{C}_t$ encapsulates the full evaluation history, enabling in-context learning to synthesise insights from prior trials and propose progressively more effective candidate designs. The regret analysis in \S~\ref{sec:theory-main} decomposes the cost of search into a surrogate-estimation term and a reservoir-quality term, with the latter shrinking as more inference-time compute is spent.

Scientific discovery tasks are further constrained by costly evaluation procedures and limited computational budget. This demand motivates \emph{model-based} bandit variants that explicitly extract knowledge from evaluation history and make the most of such knowledge in decision making. To this end, \emph{Bayesian optimisation} (BO) \citep{mockusBayesMethodsSeeking1975,jonesEfficientGlobalOptimization1998,garnettBayesianOptimization2023} provides a powerful framework for such a variant that uses a probabilistic surrogate following the rigour of Bayesian inference, where typically a Gaussian process (GP) \citep{rasmussen2006gaussian} serves as the surrogate. The GP provides predictive mean and variance, which are combined into an \emph{acquisition function} that guides the selection of the next evaluation. For example, canonical implementations such as Efficient Global Optimisation (EGO) \citep{jonesEfficientGlobalOptimization1998} and GP-UCB \citep{srinivas2010gaussian} use expected improvement (EI) and UCB \citep{auerUsingConfidenceBounds2002} as their acquisition function, respectively. Theoretical analysis shows that GP-UCB enjoys sublinear cumulative regret in terms of the maximum information gain, a measure of how informative the observations are about the objective \citep{srinivas2010gaussian}. These theoretical guarantees make BO a principled tool for experimental design for scientific discovery \citep{yuEfficientPrincipledScientific2026}.

BO has been successfully applied across chemistry, materials science, and biology, often requiring domain-specific adaptations. Examples include the PHOENICS \citep{hasePhoenicsBayesianOptimizer2018} and Gryffin \citep{haseGryffinAlgorithmBayesian2021} frameworks for chemical discovery. Combinatorial BO frameworks such as AntBO \citep{khan2022antbo} have been developed to handle the discrete sequence space of CDRH3 loops for antibody design. Benchmarking studies across multiple materials domains \citep{liangBenchmarkingPerformanceBayesian2021} have compared different surrogate models and acquisition functions, while applications in chemical synthesis \citep{shieldsBayesianReactionOptimization2021} and material discovery \citep{shoyebraihanAcceleratingMaterialDiscovery2024} further demonstrate the versatility of BO. For a comprehensive overview of BO for chemical problems, see \citep{wuRaceBottomBayesian2024}. To search among discrete sequences (e.g., SMILES and proteins), methods such as BOSS (Bayesian Optimisation over String Spaces) \citep{mossBOSSBayesianOptimization2020} use string kernels to directly construct GPs over the string space. To handle molecule design, latent-space approaches combine variational autoencoders with BO, though they struggle with validity and out-of-distribution generation \citep{griffiths2020constrained}. More recent frameworks like COMBO \citep{dreczkowskiFrameworkBenchmarksCombinatorial2023} address combinatorial and mixed-variable optimisation. Multi-objective extensions (Expected Hypervolume Improvement) \citep{ponweiserMultiobjectiveOptimizationLimited2008, daulton2021parallel} handle the competing objectives common in drug discovery and materials science.

Despite these advances, pure statistical methods face fundamental challenges. First, they struggle to propose valid candidates in vast, unbounded spaces of practical scientific problems. Next, they cannot easily incorporate high-level semantic domain knowledge. Additionally, they are often unresponsive to user-specified design constraints or preferences. These limitations create a natural opportunity for complementarity with generative models, which can propose semantically meaningful candidates and be steered by domain knowledge. This brings us to the growing body of hybrid approaches that combine the strengths of LLMs and BO.

\subsection{Hybrid Approaches: LLM-Guided Discovery}

The recognition that LLMs and statistical learning offer complementary strengths has led to a growing body of hybrid methods. A high-level approach is OPRO \citep{yangLargeLanguageModels2024}, which uses an LLM as an optimizer by describing the task in natural language and iteratively generating solutions from a prompt containing a history of solution--score pairs. While powerful in prompt optimisation, OPRO lacks the calibrated uncertainty and principled exploration of a Bayesian surrogate.

Several works have integrated LLMs directly into the BO pipeline. LLAMBO \citep{liuLargeLanguageModels2024} uses LLMs for zero-shot warm-starting, surrogate modelling, and candidate sampling, showing gains in low-data hyperparameter optimisation regimes. BoChemian \citep{rankovicBoChemianLargeLanguage2023} and its extension \citep{rankovicLargeLanguageModels2025a} use LLM embeddings as features for BO over chemical reactions, demonstrating that fine-tuning LLMs with uncertainty-aware objectives can improve discovery rates. LABO \citep{chenLABOLLMAcceleratedBayesian2026} proposes a gating criterion to dynamically balance LLM predictions against experimental observations. Other works use LLMs to generate or refine kernels for GP surrogates \citep{suwandiAdaptiveKernelDesign2025}, handle natural language feedback \citep{kobalczykLILOBayesianOptimization2025}, or perform analog circuit design \citep{yinADOLLMAnalogDesign2024,chenLLMEnhancedBayesianOptimization2024}. The LLM-as-warm-start paradigm, where the LLM suggests the initial candidates or a search region, is also explored in \citep{ramosBayesianOptimizationCatalysis2025, yinADOLLMAnalogDesign2024}. In catalysis, \citet{ramosBayesianOptimizationCatalysis2025} use frozen LLMs for in-context learning to enable BO without feature engineering, demonstrating the potential of LLMs to operate directly in language space.

Despite various attempts to combine the advantages of LLMs and BO, some recent studies have challenged their solidity. A sobering perspective is provided by \citet{guptaLLMsBayesianOptimization2025}, who find that current LLMs show no sensitivity to experimental feedback in BO tasks, and classical methods such as linear bandits and GP optimisation consistently outperform LLM agents. Similarly, \citet{kristiadiSoberLookLLMs2024} take a dispassionate stance, concluding that uncertainty taken directly from point-estimated LLMs is unreliable, and that LLMs are useful for BO over molecules primarily only when they serve as fixed feature extractors for principled GP surrogates and are pretrained or finetuned on domain-specific data. These critiques directly motivate LDM's design: we retain an explicit calibrated posterior (the GP surrogate) and use the LLM only as a candidate reservoir.

Another closely related work to LDM is dLLM \citep{yuan2026diffusion}, which for \emph{offline} black-box optimisation runs a masked-diffusion tree search where leaf candidates are scored by expected improvement under a GP fitted to an offline dataset. In contrast, LDM targets the standard \emph{online} recurrent loop with sequential, expensive evaluations. The core is that LDM must raise designs that are valuable not only for the current trial but also for gaining knowledge for the future.

In summary, LLMs provide powerful generative priors but lack calibration and robust optimisation; meanwhile, statistical methods provide principled, uncertainty-aware search but struggle to propose candidates in semantic, combinatorial spaces. LDM bridges this gap by making the acquisition function the value signal of an LLM-driven inference-time search, creating a principled, model-based framework for scientific discovery.
